\documentclass[12pt,a4paper]{article}
\usepackage[utf8x]{inputenc}
\usepackage[english,hebrew]{babel}
\usepackage{graphicx}
\usepackage{verbatim}
\usepackage{url}

\usepackage{tikz}
\usetikzlibrary{calc, external, intersections}
\tikzexternalize[prefix=tikz/]

\textwidth=15.5cm
\textheight=23cm
\topmargin=0pt
\headheight=0pt
\oddsidemargin=2em
\headsep=0pt
\renewcommand{\baselinestretch}{1.1}
\setlength{\parskip}{0.3\baselineskip plus 1pt minus 1pt}
\parindent=0pt

\newcommand*{\qd}{קוודרטריקס}

\begin{document}
\thispagestyle{empty}

\selectlanguage{hebrew}

\begin{center}
\textbf{\Huge%
איך לחלק זווית לשלושה\\
\medskip
)אם אתם מוכנים לרמות(
}

\bigskip

\bigskip

\textbf{\LARGE מוטי בן-ארי}

\bigskip

\url{http://www.weizmann.ac.il/sci-tea/benari/}

\bigskip

\selectlanguage{english}
\copyright{}\  2017 by Moti Ben-Ari.
\end{center}
\selectlanguage{english}

{\footnotesize This work is licensed under the Creative Commons Attribution-ShareAlike 3.0 Unported License. To view a copy of this license, visit \url{http://creativecommons.org/licenses/by-sa/3.0/} or send a letter to Creative Commons, 444 Castro Street, Suite 900, Mountain View, California, 94041, USA.}

\begin{center}
\includegraphics[width=.15\textwidth]{../by-sa.png}
\end{center}

\selectlanguage{hebrew}

%%%%%%%%%%%%%%%%%%%%%%%%%%%%%%%%%%%%%%%%%%%%%%%%%%%%%%%%%%%%%

\section{מבוא}

ידוע שלא ניתן לחלק זווית שרירותית לשלושה בעזרת מחוגה וסרגל. הסיבה היא שחלוקת זווית לשלושה דורשת בנייה של שורש שלישי, אבל עם מחוגה וסרגל ניתן לבנות רק אורכים המתקבלים מארבעת פעולות החשבון וכן שורש ריבועי.

המתמטיקאים היוונים גילו שבאמצעות כלים אחרים ניתן לחלק זווית לשלושה. סעיף~%
\ref{s.neusis}
מציג בנייה של ארכימדס עם כלי פשוט הנקרא ביוונית ניאוסיס
\L{(neusis)}.
סעיף~%
\ref{s.q}
מביא בנייה מסובכת יותר של היפיאס באמצעות
\qd{}
\L{(quadratrix)}.
כהטבה מיוחדת, נראה בסעיף~%
\ref{s.square}
איך ניתן לרבע מעגל באמצעות 
\qd{}.


מקורות:

\selectlanguage{english}
\noindent\url{https://en.wikipedia.org/wiki/Angle_trisection}\\
\url{https://en.wikipedia.org/wiki/Quadratrix_of_Hippias}\\
\url{https://en.wikipedia.org/wiki/Neusis_construction}
\selectlanguage{hebrew}

\newpage

%%%%%%%%%%%%%%%%%%%%%%%%%%%%%%%%%%%%%%%%%%%%%%%%%%%%%%%%%%%%%

\section{חלוקת זווית לשלושה באמצעות ניאוסיס}\label{s.neusis}

בבית הספר אנו לומדים לבנות צורות גיאומטריות באמצעות סרגל ומחוגה. למעשה, השימוש במילה "סרגל" מטעה, כי הכוונה היא למקל ישר ללא כל סימן, שהפעולה היחידה שניתן לעשות איתו היא למתוח קו ישר בין שתי נקודות. לסרגל המוכר יש סימנים המאפשרים למדוד אורכים. כדי לחלק זווית לשלושה, נשתמש ב-%
\textbf{ניאוסיס}
שהוא מקל עם שני סימנים בלבד, כאשר למען הנוחיות נגדיר שהמרחק בין שני הסימנים הוא
$1$:
\begin{center}
\selectlanguage{english}
\begin{tikzpicture}[scale=3.5]
\draw (-1,1.05) rectangle +(3.2,.1);
\draw[thick] (1.89,1.05) -- +(0,.1);
\draw[thick] (.73,1.05) -- +(0,.1);
\draw[<->] (.73,1.25) -- node[fill=white] {$1$} (1.89,1.25);
\end{tikzpicture}
\end{center}
תהי 
$\alpha$
זווית שרירותית
$\angle ABE$
בתוך מעגל שמרכזו
$B$
עם רדיוס
$1$.
ניתן לבנות את המעגל עם רדיוס
$1$
על ידי קביעת המרחק בין רגלי החוגה למרחק בין סימני הניאוסיס.

הרחיבו את רדיוס
$EB$
מחוץ למעגל. הניחו את הניאוסיס על הנקודה
$A$
והזיזו אותו עד שהוא חותך את ההרחבה של 
$EB$
בנקודה 
$D$
ואת המעגל בנקודה
$C$.
השתמשו בסימנים כדי שאורך קטע הקו
$CD$
יהיה
$1$.
ציירו את הקו
$AD$:
\begin{center}
\selectlanguage{english}
\begin{tikzpicture}[scale=3]
\coordinate (origin) at (0,0) node[below] {$B$} ;
\draw[name path=circle] (origin) circle [radius=1];
\draw (origin) node[above left,xshift=-4pt] {$\alpha$} -- (120:1) coordinate (a) node[below,xshift=-2pt] {$A$} ;
\draw (-1,0) -- (2.2,0);
\path[name path=ad] (a) -- (0,0 -| 2,0) coordinate (d) node[below] {$D$} ;
\path[name intersections={of=circle and ad,by={c,a1}}];
\fill (origin) circle [radius=.5pt];
\fill (a) circle [radius=.5pt];
\fill (c) circle [radius=.5pt] node [below,xshift=-4pt] {$C$};
\fill (d) circle [radius=.5pt];
\fill (-1,0) circle [radius=.5pt] node [left] {$E$};
\begin{scope}[rotate=-19,yshift=-11.25pt]
\draw (-1,1.05) rectangle +(3.2,.1);
\draw[thick] (1.89,1.05) -- +(0,.1);
\draw[thick] (.76,1.05) -- +(0,.1);
\draw[<->] (.73,1.25) -- node[fill=white] {$1$} (1.89,1.25);
\end{scope}
\end{tikzpicture}
\end{center}
ציירו את הקו 
$BC$
וסמנו את הזוויות וקטע הקו כפי שמופיע בתרשים שלהלן:
\begin{center}
\selectlanguage{english}
\begin{tikzpicture}[scale=3]
\coordinate (origin) at (0,0) node[below] {$B$} ;
\draw[name path=circle] (origin) circle [radius=1];
\draw (origin) node[above left,xshift=-4pt] {$\alpha$} node[above,xshift=4pt,yshift=2pt] {$\delta$} node[above right,xshift=44pt,yshift=-2pt] {$\beta$} -- node[fill=white] {$1$} (120:1) coordinate (a) node[above left] {$A$} ;
\draw (-1,0) -- (2.2,0);
\draw[name path=ad] (a) node[below right,xshift=8pt,yshift=-6pt] {$\gamma$} -- (0,0 -| 2,0) coordinate (d) node[left,xshift=-50pt,yshift=7pt] {$\beta$} node[above right] {$D$} ;
\path[name intersections={of=circle and ad,by={c,a1}}];
\draw (origin) -- node[fill=white] {$1$}(c) node[above right] {$C$} node[left,xshift=-12pt] {$\gamma$} node[below,xshift=-2pt,yshift=-4pt] {$\epsilon$};
\fill (origin) circle [radius=.5pt];
\fill (a) circle [radius=.5pt];
\fill (d) circle [radius=.5pt];
\fill (c) circle [radius=.5pt];
\fill (-1,0) circle [radius=.5pt];
\path (c) -- node[fill=white] {$1$} (d);
\end{tikzpicture}
\end{center}
השתמשנו בעובדה ש-%
$\triangle ABC$
ו-%
$\triangle BCD$
משולשים שווי שוקיים:
$AB=BC$
כי שניהם רדיוסים ו-%
$BC=CD$
לפי הבניה באמצעות הניאוסיס.

החישוב שלהלן משתמש בעובדות שסכום הזוויות של משולש ושל זוויות משלימות הוא
$\pi$
רדיאנים.
\begin{eqnarray*}
\epsilon &=& \pi - 2\beta\\
\gamma &=& \pi - \epsilon = 2\beta\\
\delta &=& \pi - 2\gamma = \pi - 4\beta\\
\alpha &=& \pi - \delta - \beta\\
&=& 4\beta -\beta\\
&=& 3\beta\,.
\end{eqnarray*}
מכאן שהזווית
$\beta$
היא שליש הזווית
$\alpha$.
%%%%%%%%%%%%%%%%%%%%%%%%%%%%%%%%%%%%%%%%%%%%%%%%%%%%%%%%%%%%%

\section{חלוקת זווית לשלושה באמצעות
\qd{}%
}\label{s.q}

התרשים שלהלן מראה
\textbf{מחוגת \qd{}}:
\begin{center}
\selectlanguage{english}
\begin{tikzpicture}[scale=.8,domain=.03:1.555,samples=100]
\draw (.1,.2) node[below left,xshift=-8pt] {$A$} -- (7.8,.2) node[below right] {$B$} -- (8,7.8) node[above right] {$C$} -- (.1,7.8) node[above left,xshift=-8pt] {$D$} -- cycle;
%\draw (.1,7.8) -- (.1,.2) -- (8,.1) -- (8,7.8);
\draw[rounded corners,rotate=60] (0,-.2) rectangle (8.2,.2);
\draw[rounded corners] (-.1,4.8) rectangle (8.2,5.2);
\fill (2.9,5) circle [radius=3pt];
\draw[thick] plot (4.6*.637*\x,{12.2*.637*\x*cot(\x r)});
\end{tikzpicture}
\end{center}

מחוגת
\qd{}
מורכב משני סרגלים )ללא סימנים כמקובל בבנייה גיאומטרית( המחוברים במפרק המאלץ אותם לנוע ביחד. סרגל אחד חייב לנוע במקביל לציר ה-%
$x$
מ-%
$DC$
עד
$AB$,
בעוד הסרגל השני מחובר במפרק למרכז מערכת הצירים ב-%
$A$.
סרגל זה מונח לאורך ציר ה-%
$y$, $AD$,
והוא מסתובב עד שהוא מונח לאורך ציר ה-%
$x$, $AB$.
העקומה המצוירת על ידי המפרק המחבר את שני הסרגלים נקראת
\textbf{עקומת ה%
\qd{}%
}
או פשוט
\textbf{\qd{}}.

כאשר מזיזים את הסרגל האופקי במהירות אחידה, החיבור עם הסרגל השני מאלץ אותו להסתובב במהירות זוותית קבועה. למעשה זו ההגדרה של ה%
\qd{}.
כאשר קואורדינטת ה-%
$y$
של הסרגל האופקי יורד מ-%
$1$
ל-%
$0$,
הזווית של הסרגל השני יחסית לציר ה-%
$x$
יורד מ-%
$\pi/2$
ל-%
$0$.

התרשים שלהלן מראה איך אפשר לחלוק זווית שרירותית 
$\alpha$
לשלושה באמצעות
\qd{}:
\begin{center}
\selectlanguage{english}
\begin{tikzpicture}[scale=.8,domain=.03:1.562,samples=100]
\draw (.1,7.8) coordinate (start) node[above left] {$D$} node[below right,xshift=32pt] {$\theta$} -- (.1,.2) node[below left] {$A$} -- (8,.1) node[below right] {$B$} -- (8,7.8) node[above right] {$C$} -- cycle;
\draw[name path=curve,thick] plot (4.6*.637*\x,{12.2*.637*\x*cot(\x r)});
\path[name path=twenty] (start) -- +(-35:9);
\path[name path=sixty] (start) -- +(-50:9);
\path[name path=xaxis] (.1,.2) -- (8,.1);
\path[name intersections={of=twenty and curve,by={x1,tri}}];
\draw (start) -- (tri);
\fill (tri) circle [radius=2pt] node[right] {$P_2$};
\path[name intersections={of=sixty and curve,by={x2,angle}}];
\fill (angle) circle [radius=2pt] node[right] {$P_1$};
\draw (start) -- (angle);
\path[name intersections={of=xaxis and curve,by=x}];
\fill (x) circle [radius=2pt];
%\path (.1,.2) -- node[below] {$2/\pi$} (x);
\draw[dashed] (tri) -- (tri -| .1,.2) coordinate (t3);
\fill (t3) circle [radius=2pt] node[left] {$F$};
\draw[dashed] (angle) -- (angle -| .1,.2) coordinate (t);
\fill (t) circle [radius=2pt] node[left] {$E$};
\draw[<->] (-1.2,7.8) -- node[fill=white] {$t/3$} (-1.2,7.8 |- t3);
\draw[<->] (-.6,7.8) -- node[fill=white] {$t$} (-.6,7.8 |- t);
\draw[<->] (-.6,7.8 |- t) -- node[fill=white] {$1-t$} (-.6,.2);
\draw (3.5,7.8) arc[start angle=0,delta angle=-49,radius=3.5];
\draw (1,7.8) arc[start angle=0,delta angle=-32,radius=1];
\node at (3.3,6) {$\alpha$}; 
\end{tikzpicture}
\end{center}
הזווית
$\alpha$
היא
$\angle CDP_1$,
כאשר הנקודה
$P_1$
היא החיתוך בין הקו המגדיר את הזווית 
$\alpha$
לבין ה%
\qd{}.
קואורדינטת ה-%
$y$
שלה היא
$1-t$,
כאשר 
$t$
הוא המרחק שהסרגל האופקי נע ממקומו ההתחלתי
$DC$.
חלקו את
\textbf{קטע הקו}
$DE$
לשלושה חלקים כדי לקבל את הנקודה
$F$.
יהי
$P_2$
נקודת החיתוך בין הקו מ-%
$F$
המקביל ל-%
$DC$
לבין ה%
\qd{}.
לפי העיקרון של מהירויות שוות:
\begin{eqnarray*}
\frac{\theta}{\alpha} &=& \frac{t/3}{t}\\
&&\\
\theta &=& \alpha/3\,.
\end{eqnarray*}

\newpage

%%%%%%%%%%%%%%%%%%%%%%%%%%%%%%%%%%%%%%%%%%%%%%%%%%%%%%%%%%%%%

\section{ריבוע המעגל באמצעות
\qd{}%
}\label{s.square}

נחשב את המשוואה של ה%
\qd{}
לפי התרשים שלהלו:
\begin{center}
\selectlanguage{english}
\begin{tikzpicture}[scale=.8,domain=.01:1.57,samples=100]
\draw (0,0) node[below left] {$A$} node [above right,xshift=10pt] {$\theta$} -- (8,0) node[below right] {$B$} -- (8,8) node[above right] {$C$} -- (0,8) node[above left] {$D$} -- cycle;
\draw[name path=horiz] (0,5) -- (8,5);
\draw[name path=slant] (0,0) -- (61:8);
\path[name intersections={of=horiz and slant,by=joint}];
\draw[dashed] (joint) -- node[right] {$y$} (joint |- 0,0) coordinate (f);
\path (0,0) -- node[below] {$x$} (0,0 -| joint) node[below] {$F$};
\path (0,5) node[left] {$E$} -- node[left] {$t$} (0,8);
\path (0,0) -- node[left] {$1-t$} (0,5);
\fill (joint) circle [radius=2pt] node[above right,xshift=4pt] {$P$};
\fill (f) circle [radius=2pt];
\fill (0,5) circle [radius=2pt];
\fill (4.28,0) circle [radius=2pt] node[below] {$G$};
\draw[name path=curve,thick] plot (4.3*.637*\x,{12.5*.637*\x*cot(\x r)});
\end{tikzpicture}
\end{center}
נניח שהסרגל האופקי נע מרחק 
$t$
לאורך ציר ה-%
$y$
עד לנקודה
$E$,
והסרגל המסתובב מגדיר זווית 
$\theta$
עם ציר ה-%
$x$.
הנקודה 
$P$
היא החיתוך בין
\qd{}
לבין הסרגל האופקי, והנקודה
$F$
היא היטל של 
$P$
על ציר ה-%
$x$.
מהן הקואורדינטות של הנקודה 
$P$
על ה%
\qd{}?
ברור ש:
\[
y=PF=EA=1-t\,.
\]
על העקומה,
$\theta$
יורד באותו קצב ש-%
$t$
עולה:
\begin{eqnarray*}
\frac{1-t}{1} &=& \frac{\theta}{\pi/2}\\
&&\\
\theta &=&\frac{\pi}{2}(1-t)\,.
\end{eqnarray*}
נבדוק אם זה הגיוני: כאשר
$t=0$
אז
$\theta=\pi/2$,
וכאשר
$t=1$
אז
$\theta=0$.

את קואורדינטת ה-%
$x$
של
$P$
נקבל בטריגונומטריה:
\[
\tan \theta = \frac{y}{x}\,.
\]
ומכאן:
\[
x = \frac{y}{\tan\theta}=y\cot\theta=y\cot \frac{\pi}{2}(1-t)=y\cot \frac{\pi}{2}y\,.
\]
בדרך כלל המשוואה של עקומה היא מהצורה
$y=f(x)$,
אבל אפשר גם להשתמש במשוואה מהצורה
$x=f(y)$.

נחשב את קוארדינטת ה-%
$x$
של הנקודה 
$G$,
החיתוך של ה%
\qd{}
עם ציר ה-%
$x$.
לא ניתן פשוט להציב
$y=0$
כי
$\cot 0$
לא מוגדר, אבל ייתכן שיהיה לנו מזל אם נחשב את הגבול של
$x$
כאשר 
$y$
שואף ל-%
$0$:
\[
x = y\cot \frac{\pi}{2}y = \frac{2}{\pi}\cdot \frac{\pi}{2}y\cot \frac{\pi}{2}y\,.
\]
למען הנוחיות, נחליף משתנה
$z=\frac{\pi}{2}y$,
ואז נחשב את הגבול:
\[
\lim_{z\rightarrow 0} z\cot z = \lim_{z\rightarrow 0} \frac{z\cos z}{\sin z} = \lim_{z\rightarrow 0} \frac{\cos z}{\frac{\sin z}{z}} = \frac{\cos 0}{1} = 1\,,
\]
השתמשנו בעובדה הידועה ש-%
$\displaystyle\lim_{z\rightarrow 0} \frac{\sin z}{z}=1$.

כאשר
$y\rightarrow 0$:
\[
x \rightarrow \frac{2}{\pi}\cdot \lim_{y\rightarrow 0}\frac{\pi}{2}y\cot \frac{\pi}{2}y = \frac{2}{\pi}\cdot 1 = \frac{2}{\pi}\,.
\]
על ידי שימוש ב%
\qd{}
בנינו קטע קו
$AG$
שאורכו
$x=\displaystyle\frac{2}{\pi}$.
עם סרגל רגיל ומחוגה, קל לבנות קו באורך
$\sqrt{\pi}=\sqrt{\frac{2}{x}}$,
ואז לבנות ריבוע ששטחו
$\pi$.


\end{document}
